%! Parametric Functions Modeling Planar Motion — Unit 4, Lesson 4.2 — Exit Ticket (Answer Key)
\documentclass[10pt]{article}
\usepackage{precalculus-article}
\usepackage{precalculus-key}   % pulls in -boxes; do NOT also load -boxes
\newcommand{\TallMath}[1]{$\displaystyle #1\rule[-1.4em]{0pt}{3.2em}$}

\begin{document}

\pageheader{Unit 4, Lesson 4.2}{Exit Ticket --- Answer Key}
\namedateperiod

\vspace{0.10in}

\begin{notesbox}{Planar Motion Check --- Key}

A particle moves with $f(t) = (t^2 - 1,\; 2t - t^2)$ for $0 \le t \le 3$.

\vspace{8pt}
\textbf{1.}\quad Complete the table.

\vspace{4pt}
\begin{center}
\renewcommand{\arraystretch}{1.7}
\begin{tabular}{|c|c|c|c|}
\hline
$t$ & $x(t) = t^2 - 1$ & $y(t) = 2t - t^2$ & $(x,\; y)$ \\
\hline
$0$ & \ans{$-1$} & \ans{$0$}  & \ans{$(-1,\;0)$}  \\
\hline
$1$ & \ans{$0$}  & \ans{$1$}  & \ans{$(0,\;1)$}   \\
\hline
$2$ & \ans{$3$}  & \ans{$0$}  & \ans{$(3,\;0)$}  \\
\hline
$3$ & \ans{$8$}  & \ans{$-3$} & \ans{$(8,\;-3)$}  \\
\hline
\end{tabular}
\end{center}

\vspace{10pt}
\textbf{2.}\quad Horizontal extrema:

\ans{$x(t) = t^2 - 1$ is an upward parabola. Vertex at $t = 0$.\\
Minimum: $x(0) = -1$ (leftmost position).\\
Maximum on $[0,3]$: $x(3) = 8$ (rightmost position).}

\vspace{0.06in}
\textbf{3.}\quad When does the particle cross the $y$-axis ($x = 0$)?

\ans{$t^2 - 1 = 0 \;\Rightarrow\; (t-1)(t+1) = 0 \;\Rightarrow\;
t = 1$ (only $t = 1$ is in $[0,3]$).\\
At $t = 1$: $(x(1), y(1)) = (0,\, 1)$. The particle crosses the
$y$-axis at $(0,\,1)$.}

\begin{teachernote}
  Q2: Students who report $t = 0$ as the minimum should be told to
  evaluate $x(0) = -1$ --- the minimum \emph{value} is $-1$, not the
  time. Q3: $t = -1$ is not in the domain; accept only $t = 1$.
  Students who report $(0, 0)$ confused $y(0) = 0$ with the intercept
  coordinates --- redirect them to evaluate $y(1) = 1$.
\end{teachernote}

\end{notesbox}

\end{document}
